\section*{Ethics Statement}
In conducting this research, we placed a strong emphasis on ethical considerations, particularly due to the sensitive nature of legal data and the methodologies employed. The \texttt{\named} dataset, used extensively in this study, was sourced from publicly accessible legal search engines, ensuring compliance with data privacy and usage regulations. We have taken steps to remove any meta-information such as judge names, case titles, and case IDs to protect the privacy and confidentiality of the individuals involved.

Furthermore, the computational resources utilized in this study were obtained through ethical and legitimate means. We subscribed to Google Colab Pro and other necessary cloud services, ensuring that all resources used for model training and testing were accessed legally. This financial support not only facilitated our research but also contributed to the sustainability of these services.

In addition to adhering to legal and ethical guidelines in data handling and resource usage, we are committed to transparency and reproducibility in our research. The \texttt{\named} dataset and the code for our models, including \texttt{\namel}, for now, has been made available to promote open science and enable other researchers to replicate and build upon our work.

Finally, we acknowledge the potential societal impact of deploying AI in legal settings. Our models are designed to assist, not replace, human judgment, and we stress the importance of human oversight in any AI-assisted legal decision-making process. We remain committed to ongoing ethical scrutiny as we advance this research field.